\documentclass[11pt,a4paper]{article}
% =====================================================================
%  Gemeinsame Vorlage der Arbeitsblätter "Graphen" – Informatik 11 (NTG)
%  (gleiches Layout wie informatik/10/java/arbeitsblaetter/ab-vorlage.tex)
%  Einbinden in jedem Blatt direkt nach \documentclass: % =====================================================================
%  Gemeinsame Vorlage der Arbeitsblätter "Graphen" – Informatik 11 (NTG)
%  (gleiches Layout wie informatik/10/java/arbeitsblaetter/ab-vorlage.tex)
%  Einbinden in jedem Blatt direkt nach \documentclass: % =====================================================================
%  Gemeinsame Vorlage der Arbeitsblätter "Graphen" – Informatik 11 (NTG)
%  (gleiches Layout wie informatik/10/java/arbeitsblaetter/ab-vorlage.tex)
%  Einbinden in jedem Blatt direkt nach \documentclass: \input{ab-vorlage}
%  Übersetzen mit XeLaTeX, LuaLaTeX, pdfLaTeX oder Tectonic.
%
%  Blatt:   xelatex ab-3-1-grundbegriffe.tex
%  Lösung:  xelatex -jobname=ab-3-1-grundbegriffe-lsg "\def\mitloesung{}\input{ab-3-1-grundbegriffe}"
%           (füllt alle Lücken rot aus; siehe \lsg, \lsglinien, \lsgoder)
% =====================================================================
\usepackage[a4paper,left=18mm,right=18mm,top=26mm,bottom=17mm,headheight=15mm,headsep=4mm,footskip=8mm]{geometry}
\usepackage{iftex}
\ifPDFTeX
  \usepackage[T1]{fontenc}
  \usepackage[utf8]{inputenc}
  \usepackage{helvet}
  \usepackage[scaled=0.86]{beramono}
  \renewcommand{\familydefault}{\sfdefault}
\else
  \usepackage{fontspec}
  \setmainfont{texgyreheros}[Extension=.otf,UprightFont=*-regular,BoldFont=*-bold,ItalicFont=*-italic,BoldItalicFont=*-bolditalic]
  \setmonofont{DejaVuSansMono}[Extension=.ttf,UprightFont=*,BoldFont=*-Bold,ItalicFont=*-Oblique,BoldItalicFont=*-BoldOblique,Scale=0.86]
\fi
\usepackage[ngerman,shorthands=off]{babel}
\usepackage{xcolor,graphicx,tabularx,array,enumitem,fancyhdr,tikz,listings,multicol,calc,etoolbox}
\usepackage[most]{tcolorbox}
\usetikzlibrary{arrows.meta,positioning,calc,decorations.pathreplacing,shapes.geometric}
\usepackage[hidelinks]{hyperref}
\setlength{\parindent}{0pt}
\setlength{\parskip}{3pt}
\setlist{nosep,leftmargin=*}

% ---------- Lösungsschalter ----------
\newif\ifloesung
\ifdefined\mitloesung\loesungtrue\fi
\newcommand{\lsgfarbe}{\color{red!75!black}}

% ---------- Kopf- und Fußzeile (einheitliche Form) ----------
\newcommand{\lehrkraft}{Hau}
\newcommand{\themenbereich}{3. Graphen}
\newcommand{\abnummer}{}
\pagestyle{fancy}\fancyhf{}
\renewcommand{\headrulewidth}{0.4pt}
\fancyhead[L]{\small Klasse 11\\Datum:}
\fancyhead[C]{\small Themenbereich:\\\themenbereich}
\fancyhead[R]{\small \lehrkraft\ifloesung\\{\lsgfarbe\bfseries Lösung}\fi}
\fancyfoot[C]{\scriptsize edu-mrh.de\,·\,Informatik 11\,·\,Blatt \abnummer\ifloesung\ (Lösung)\fi\,·\,Seite \thepage}
\newcommand{\abtitel}[2]{\renewcommand{\abnummer}{#1}\begin{center}\Large\bfseries #1\quad\underline{#2}\end{center}\vspace{-1mm}}

% ---------- Boxen ----------
\newenvironment{aufgabe}[2][]{\begin{tcolorbox}[enhanced,breakable,colback=white,colframe=black,boxrule=0.7pt,arc=3mm,left=2mm,right=2mm,top=1.5mm,bottom=1.5mm,before skip=2.5mm,after skip=2.5mm]\textbf{Aufgabe #2)}\ifstrempty{#1}{}{ \textit{#1}}\par\vspace{0.6mm}}{\end{tcolorbox}}
\newtcolorbox{merke}[1][Merke]{enhanced,breakable,colback=green!5,colframe=green!40!black,boxrule=0.8pt,arc=2mm,left=2mm,right=2mm,top=1.5mm,bottom=1.5mm,before skip=2.5mm,after skip=2.5mm,title={#1},fonttitle=\bfseries,coltitle=white}
\newtcolorbox{hinweis}{enhanced,breakable,colback=yellow!12,colframe=orange!70!black,boxrule=0.6pt,arc=2mm,left=2mm,right=2mm,top=1mm,bottom=1mm,before skip=2mm,after skip=2mm}
\newtcolorbox{ausblick}[1][Ausblick]{enhanced,breakable,colback=teal!5,colframe=teal!60!black,boxrule=0.8pt,arc=2mm,left=2mm,right=2mm,top=1.5mm,bottom=1.5mm,before skip=2.5mm,after skip=2.5mm,title={#1},fonttitle=\bfseries,coltitle=white}
\newcommand{\web}[1]{{\small\color{gray}\textit{Website: #1}}}

% ---------- Lücken, Linien, Karos ----------
\newcommand{\luecke}[1][3.5cm]{\rule[-0.5ex]{#1}{0.4pt}}
\newcommand{\linien}[1]{\par\vspace{1.5mm}\foreach \i in {1,...,#1}{\noindent\rule[0pt]{\linewidth}{0.3pt}\par\vspace{5.5mm}}\vspace{-3mm}}
\newcommand{\kariert}[1]{\par\vspace{1.5mm}\noindent\begin{tikzpicture}\draw[step=5mm,gray!55,thin] (0,0) grid (\linewidth-1pt,#1*5mm);\end{tikzpicture}\par}
\newcommand{\karobox}[2]{\begin{tikzpicture}\draw[step=5mm,gray!55,thin] (0,0) grid (#1,#2);\end{tikzpicture}}
\newcommand{\code}[1]{\texttt{#1}}

% Lücke mit Lösung: \lsg[Mindestbreite]{Lösungstext}
\newlength{\lsgbreite}
\newcommand{\lsg}[2][3.5cm]{\ifloesung\settowidth{\lsgbreite}{\,#2\,}\ifdim\lsgbreite<#1\setlength{\lsgbreite}{#1}\fi\underline{\makebox[\lsgbreite][l]{\lsgfarbe\,#2}}\else\luecke[#1]\fi}
% Schreiblinien mit Lösung: \lsglinien{Anzahl}{Lösungstext} (gleiche Höhe wie \linien)
\newcommand{\lsglinien}[2]{\ifloesung\par\vspace{1.5mm}\noindent\begin{minipage}[t][#1\dimexpr 6.9mm\relax][t]{\linewidth}\lsgfarbe #2\end{minipage}\par\vspace{-1.5mm}\else\linien{#1}\fi}
% Beliebiger Inhalt: \lsgoder{nur in der Lösung}{nur im Blatt}
\newcommand{\lsgoder}[2]{\ifloesung#1\else#2\fi}

% ---------- Graphen zeichnen ----------
\tikzset{
  knoten/.style={draw,rounded corners=3pt,fill=orange!12,minimum height=6mm,inner sep=3pt,font=\small},
  kknoten/.style={draw,circle,fill=orange!12,minimum size=7mm,inner sep=1pt,font=\small},
  kante/.style={thick},
  gkante/.style={thick,-{Stealth[length=2.6mm]}},
  dkante/.style={thick,{Stealth[length=2.6mm]}-{Stealth[length=2.6mm]}},
  gewicht/.style={font=\scriptsize,fill=white,inner sep=1pt,text=teal!60!black},
  lsgkante/.style={thick,red!75!black},
  lsggkante/.style={thick,red!75!black,-{Stealth[length=2.6mm]}},
  lsgknoten/.style={knoten,draw=red!75!black,text=red!75!black,fill=red!4},
}
% Adjazenzmatrix zum Ausfüllen: \adjmx[Lösungseinträge]{Anzahl}{Knotenliste}{Zellbreite}
% Lösungseinträge als Zeile/Spalte/Text, z. B. \adjmx[1/2/X,2/1/X]{3}{A,B,C}{7mm}
\newcommand{\adjmx}[4][]{\begin{tikzpicture}[x=#4,y=#4]
  \foreach \n [count=\i] in {#3} {\node[font=\small\bfseries] at (\i,0) {\n}; \node[font=\small\bfseries,anchor=east] at (0.45,-\i) {\n};}
  \foreach \k in {0,...,#2} {\draw[gray] (\k+0.5,-0.5) -- (\k+0.5,-#2-0.5); \draw[gray] (0.5,-\k-0.5) -- (#2+0.5,-\k-0.5);}
  \ifloesung\ifstrempty{#1}{}{\foreach \z/\s/\t in {#1} {\node[text=red!75!black,font=\small] at (\s,-\z) {\t};}}\fi
\end{tikzpicture}}

% ---------- Java-Listings ----------
\lstset{language=Java,basicstyle=\ttfamily\small,keywordstyle=\bfseries,commentstyle=\color{gray!80!black},stringstyle=\color{black},showstringspaces=false,columns=flexible,keepspaces=true,tabsize=3,frame=single,framerule=0.4pt,rulecolor=\color{black},xleftmargin=2mm,framexleftmargin=1mm,aboveskip=2mm,belowskip=2mm,breaklines=true,escapechar=§,
 literate={ä}{{\"a}}1 {ö}{{\"o}}1 {ü}{{\"u}}1 {Ä}{{\"A}}1 {Ö}{{\"O}}1 {Ü}{{\"U}}1 {ß}{{\ss}}1 {–}{{--}}1 {„}{{\glqq}}1 {“}{{\grqq}}1}
\lstnewenvironment{java}[1][]{\lstset{#1}}{}

%  Übersetzen mit XeLaTeX, LuaLaTeX, pdfLaTeX oder Tectonic.
%
%  Blatt:   xelatex ab-3-1-grundbegriffe.tex
%  Lösung:  xelatex -jobname=ab-3-1-grundbegriffe-lsg "\def\mitloesung{}\documentclass[11pt,a4paper]{article}
\input{ab-vorlage}
% Dörfer-Graph (Merkkasten Begriffe): Koordinaten wie im Begriffe-Explorer der Website
\newcommand{\doerferknoten}{%
  \node[knoten,font=\tiny,minimum height=4mm,inner sep=1.5pt] (A) at (4.4,3.3) {Altdorf};
  \node[knoten,font=\tiny,minimum height=4mm,inner sep=1.5pt] (W) at (1.2,2.9) {Weiler};
  \node[knoten,font=\tiny,minimum height=4mm,inner sep=1.5pt] (F) at (5.1,2.0) {Fischbach};
  \node[knoten,font=\tiny,minimum height=4mm,inner sep=1.5pt] (Z) at (2.9,2.0) {Ziegelstein};
  \node[knoten,font=\tiny,minimum height=4mm,inner sep=1.5pt] (N) at (0.9,0.95) {Neustadt};
  \node[knoten,font=\tiny,minimum height=4mm,inner sep=1.5pt] (B) at (4.4,0.8) {Burg};
  \node[knoten,font=\tiny,minimum height=4mm,inner sep=1.5pt] (R) at (2.6,0.2) {Rain};}
\begin{document}
\abtitel{3.1}{Graphen -- Grundbegriffe}

Hänsel und Gretel haben die Hexe überlistet. Doch die Brotkrumen, mit denen sie ihren Weg markiert hatten, haben die Vögel gefressen. Findest du trotzdem nach Hause?

\begin{aufgabe}[Hänsel und Gretel]{1}
Starte das Spiel auf der Website (Station 3.1).
\textbf{a)} Zeichne beim Spielen eine Karte des Waldes: jeder Ort ein Kästchen, jeder Weg eine Linie. Führt ein Weg nur in eine Richtung, zeichne einen Pfeil. Das Hexenhaus ist schon eingezeichnet.
\end{aufgabe}
\noindent\begin{tikzpicture}
\draw[step=5mm,gray!55,thin] (0,0) grid (17.4,7);
\node[knoten,fill=green!15] (hex) at (4,2.3) {Hexenhaus (Start)};
\ifloesung
  \node[lsgknoten] (warn) at (5.5,6.5) {Warnung};
  \node[lsgknoten] (wolf) at (10.5,6.5) {Werwolf};
  \node[lsgknoten] (stadt) at (1.6,5.2) {Stadt};
  \node[lsgknoten] (abz) at (5.5,5.2) {Abzweigung};
  \node[lsgknoten] (spinne) at (13,5.2) {Spinne};
  \node[lsgknoten] (baeck) at (1.6,3.6) {Bäckerei};
  \node[lsgknoten] (pilze) at (4,3.9) {Pilze};
  \node[lsgknoten] (w3) at (13,3.4) {Wald3};
  \node[lsgknoten] (w1) at (8,1.0) {Wald1};
  \node[lsgknoten] (stein) at (12.2,2.0) {Steinkreis};
  \node[lsgknoten] (w2) at (15,0.8) {Wald2};
  \node[lsgknoten] (heim) at (16.2,2.6) {Zuhause};
  \draw[lsgkante] (stadt) -- (baeck);
  \draw[lsgkante] (stadt) -- (pilze);
  \draw[lsgkante] (pilze) -- (abz);
  \draw[lsgkante] (pilze) -- (hex);
  \draw[lsggkante] (pilze) to[bend left=50] node[right,font=\scriptsize,text=red!75!black] {Pilze essen} (hex);
  \draw[lsgkante] (abz) -- (warn);
  \draw[lsggkante] (warn) -- (wolf);
  \draw[lsgkante] (abz) -- (spinne);
  \draw[lsgkante] (spinne) -- (w3);
  \draw[lsgkante] (hex) -- (w1);
  \draw[lsgkante] (w1) -- (w3);
  \draw[lsgkante] (w1) -- (w2);
  \draw[lsgkante] (w1) -- (stein);
  \draw[lsgkante] (w2) -- (w3);
  \draw[lsgkante] (w2) -- (stein);
  \draw[lsgkante] (w3) -- (stein);
  \draw[lsggkante] (w2) -- (heim);
  \node[text=red!75!black,font=\scriptsize,align=left,anchor=north west] at (0.2,1.3) {13 Knoten, 17 Kanten\\(14 ungerichtet, 3 gerichtet)};
\fi
\end{tikzpicture}

\textbf{b)} Am Ende zeigt dir das Spiel deinen Weg. Notiere ihn und die Anzahl der Schritte:
\lsglinien{1}{individuell; kürzester Weg: Hexenhaus -- Wald1 -- Wald2 -- Zuhause (3 Schritte)}

\begin{minipage}[t]{0.6\linewidth}
\begin{merke}[Graph und Kurzschreibweise $G = (V, E)$]\raggedright
Ein Graph besteht aus \lsg[2.3cm]{Knoten} (engl. \textit{vertices}) und \lsg[2.3cm]{Kanten} (engl. \textit{edges}). Eine Kante verbindet zwei Knoten. Im Spiel sind die Orte die \lsg[2cm]{Knoten}, die Wege die \lsg[2cm]{Kanten}.\\[1mm]
Kurzschreibweise $G = (V, E)$: $V$ ist die Menge der Knoten, $E$ die Menge der Kanten.
\begin{itemize}
\item $(A, B)$: \lsg[2.5cm]{gerichtete} Kante, nur von $A$ nach $B$
\item $\{A, B\}$: \lsg[2.5cm]{ungerichtete} Kante, steht für $(A, B)$ und \lsg[1.4cm]{$(B, A)$}
\end{itemize}
\end{merke}
\end{minipage}\hfill
\begin{minipage}[t]{0.37\linewidth}
\vspace{1mm}
\centering\small \textbf{Beispiel: Wer folgt wem?}\\[1mm]
\begin{tikzpicture}[x=0.95cm,y=0.9cm]
\node[knoten] (leon) at (0.7,2.2) {Leon};
\node[knoten] (letifa) at (3.3,2.4) {Letifa};
\node[knoten] (luise) at (4.3,1.1) {Luise};
\node[knoten] (hamza) at (2.4,0.3) {Hamza};
\node[knoten] (sabine) at (0.5,0.6) {Sabine};
\draw[dkante] (leon) -- (letifa);
\draw[gkante] (letifa) -- (luise);
\draw[gkante] (leon) -- (luise);
\draw[dkante] (leon) -- (hamza);
\draw[gkante] (hamza) -- (luise);
\end{tikzpicture}
\end{minipage}\\[1mm]
{\small $V = \{\text{Hamza}, \text{Leon}, \text{Letifa}, \text{Luise}, \text{Sabine}\}$\\
$E = \{\{\text{Leon}, \text{Letifa}\}, \{\text{Leon}, \text{Hamza}\}, (\text{Letifa}, \text{Luise}), (\text{Leon}, \text{Luise}), (\text{Hamza}, \text{Luise})\}$}

\begin{aufgabe}[Hänsel und Gretel]{1}
\textbf{c)} Schreibe den Ausschnitt deiner Karte mit den Knoten Hexenhaus, Pilze, Stadt, Bäckerei und Wald1 in Kurzschreibweise auf (5 Kanten).\\[2mm]
$V = \{$\lsg[15.4cm]{Hexenhaus, Pilze, Stadt, Bäckerei, Wald1}$\}$\\[3mm]
$E = \{$\lsg[15.4cm]{\{Hexenhaus, Pilze\}, (Pilze, Hexenhaus), \{Pilze, Stadt\}, \{Stadt, Bäckerei\}, \{Hexenhaus, Wald1\}}$\}$
\end{aufgabe}

\newpage
\begin{minipage}[t]{0.71\linewidth}
\begin{merke}[Pfade, Zyklen, Richtung, Gewicht]\raggedright\small
\textbf{Pfad (Weg):} Folge von Knoten, in der aufeinanderfolgende Knoten durch eine \lsg[1.8cm]{Kante} verbunden sind; keine Kante wird doppelt benutzt.\\[0.6mm]
\textbf{Einfacher Pfad:} Pfad, in dem sich kein \lsg[1.8cm]{Knoten} wiederholt.\\
Beispiel in A: \lsg[5.6cm]{Altdorf -- Weiler -- Ziegelstein -- Neustadt}\\[0.6mm]
\textbf{Länge eines Pfades:} ungewichtet die \lsg[3.2cm]{Anzahl der Kanten}, gewichtet die \lsg[3.2cm]{Summe der Gewichte}.\\[0.6mm]
\textbf{Zyklus:} einfacher Pfad, dessen \lsg[5cm]{Start- und Endknoten gleich sind}.\\
Beispiel in A: \lsg[5.6cm]{Weiler -- Ziegelstein -- Fischbach -- Weiler}\\
Graph mit Zyklus: \emph{zyklisch}, sonst \lsg[3.4cm]{azyklisch (zyklenfrei)}.\\[0.6mm]
\textbf{Gerichtet:} Kanten nur \lsg[3cm]{in Pfeilrichtung} (Pfeil).\ \textbf{Ungerichtet:} \lsg[2.8cm]{in beide Richtungen} (Strich).\\[0.6mm]
\textbf{Stark zusammenhängend:} Von \lsg[2.4cm]{jedem Knoten} gibt es einen Pfad zu \lsg[3.2cm]{jedem anderen Knoten}.\\
\textbf{Schwach zusammenhängend:} \lsg[5.4cm]{nur ohne Beachtung der Richtungen}\\ \lsg[\linewidth]{hängt alles zusammen. B ist schwach, denn von Rain führt keine Kante weg.}\\[0.6mm]
\textbf{Gewichtet:} Jede Kante trägt eine \lsg[2.6cm]{Zahl (Gewicht)}, z.\,B. \lsg[3.6cm]{Entfernung, Fahrzeit}.
\end{merke}
\end{minipage}\hfill
\begin{minipage}[t]{0.27\linewidth}
\vspace{1mm}
\textbf{\small A} {\scriptsize ungerichtet}\\
\begin{tikzpicture}[x=0.8cm,y=0.85cm]
\doerferknoten
\draw[kante] (A)--(W); \draw[kante] (F)--(A); \draw[kante] (W)--(F); \draw[kante] (W)--(Z);
\draw[kante] (Z)--(F); \draw[kante] (Z)--(N); \draw[kante] (B)--(Z); \draw[kante] (F)--(B);
\draw[kante] (N)--(R); \draw[kante] (B)--(R);
\end{tikzpicture}\\[2mm]
\textbf{\small B} {\scriptsize gerichtet}\\
\begin{tikzpicture}[x=0.8cm,y=0.85cm]
\doerferknoten
\draw[gkante] (A)--(W); \draw[gkante] (F)--(A); \draw[gkante] (W)--(F); \draw[gkante] (W)--(Z);
\draw[gkante] (Z)--(F); \draw[gkante] (Z)--(N); \draw[gkante] (B)--(Z); \draw[gkante] (F)--(B);
\draw[gkante] (N)--(R); \draw[gkante] (B)--(R);
\end{tikzpicture}
\end{minipage}

\begin{aufgabe}{2}
{\raggedright Zeichne $G = (V, E)$ mit $V = \{A, B, C, D, E, F\}$ und\\ $E = \{(A, B), (A, C), (A, E), (B, C), (C, C), (C, D), (E, C), (F, D)\}$. Beantworte die Fragen.\par}\vspace{1mm}
\begin{minipage}[t]{0.42\linewidth}
\vspace{0pt}
\begin{tikzpicture}
\draw[step=5mm,gray!55,thin] (0,0) grid (7,3);
\ifloesung
  \foreach \n/\x/\y/\t in {a/0.6/1.5/A, b/2/2.55/B, e/2/0.45/E, c/3.4/1.5/C, d/4.9/1.5/D, f/6.4/1.5/F} {
    \node[kknoten,minimum size=6mm,draw=red!75!black,text=red!75!black] (\n) at (\x,\y) {\t};}
  \draw[lsggkante] (a)--(b); \draw[lsggkante] (a)--(c); \draw[lsggkante] (a)--(e);
  \draw[lsggkante] (b)--(c); \draw[lsggkante] (c)--(d); \draw[lsggkante] (e)--(c);
  \draw[lsggkante] (f)--(d);
  \draw[lsggkante] (c) to[out=60,in=120,looseness=6] (c);
\fi
\end{tikzpicture}
\end{minipage}\hfill
\begin{minipage}[t]{0.55\linewidth}
\vspace{0pt}\small\raggedright
Gerichtet? Woran erkennbar? \lsg[3.9cm]{ja, Paare ( ) statt \{ \}}\\[1.5mm]
Ein Zyklus: \lsg[5.7cm]{(C, C), eine Schlinge (Länge 1)}\\[1.5mm]
Stark/schwach zusammenhängend? \lsg[3.6cm]{schwach: von D}\\[1.5mm]
\lsg[0.98\linewidth]{kommt man nirgends hin.}\\[1.5mm]
Länge des längsten einfachen Pfades: \lsg[3.2cm]{3 (A--B--C--D)}
\end{minipage}
\end{aufgabe}

\begin{aufgabe}[Zurück in den Wald]{3}
Deine Karte aus Aufgabe 1: \textbf{a)} gerichtet?\ \textbf{b)} ein Zyklus?\ \textbf{c)} stark oder schwach zusammenhängend?\ \textbf{d)} „Der Weg zurück … scheint verschwunden“: Ist die Kante Hexenhaus -- Wald1 gerichtet? Prüfe im Spiel.
\lsglinien{2}{\small a) gemischt; gerichtet nur (Pilze, Hexenhaus), (Warnung, Werwolf), (Wald2, Zuhause)\quad b) Wald1 -- Wald2 -- Wald3 -- Wald1\\ c) schwach: Von Zuhause/Werwolf führt keine Kante weg.\quad d) nein: Von Wald2 aus führt ein Weg zurück.}
\end{aufgabe}

\begin{aufgabe}[Ordne die Graphen ein.]{4}
\vspace{-1mm}
\small\renewcommand{\arraystretch}{1.15}
\begin{tabularx}{\linewidth}{|>{\centering\arraybackslash}m{5.6cm}|>{\centering\arraybackslash}m{1.7cm}|>{\centering\arraybackslash}m{1.1cm}|>{\centering\arraybackslash}m{1.1cm}|>{\centering\arraybackslash}m{1.1cm}|>{\centering\arraybackslash}X|}\hline
\textbf{Graph} & \scriptsize\textbf{zusammen-\newline hängend?}\newline stark/schwach/nein & \scriptsize\textbf{ge-\newline richtet?} & \scriptsize\textbf{ge-\newline wichtet?} & \scriptsize\textbf{zyklisch?} & \scriptsize\textbf{Länge eines längsten einfachen Pfades}\\\hline
\rule{0pt}{8mm}\textbf{I}\ \begin{tikzpicture}[baseline=(k5.base),x=0.85cm,y=0.55cm]
\foreach \n/\x/\y in {1/0/1,2/1/1,3/2/1,4/3/1,5/1/0,6/3/0} {\node[kknoten,minimum size=5mm,inner sep=0pt,font=\tiny] (k\n) at (\x,\y) {K\n};}
\draw[kante] (k1)--(k2)--(k3)--(k4); \draw[kante] (k2)--(k5);
\end{tikzpicture}
& \lsgoder{\lsgfarbe nein}{} & \lsgoder{\lsgfarbe nein}{} & \lsgoder{\lsgfarbe nein}{} & \lsgoder{\lsgfarbe nein}{} & \lsgoder{\lsgfarbe 3: K1--K2--K3--K4}{}\\\hline
\rule{0pt}{5mm}\textbf{II}\ $V = \{A, B, C\}$\newline $E =\{(A, B), (B, A), (C, B), (C, C)\}$
& \lsgoder{\lsgfarbe schwach}{} & \lsgoder{\lsgfarbe ja}{} & \lsgoder{\lsgfarbe nein}{} & \lsgoder{\lsgfarbe ja}{} & \lsgoder{\lsgfarbe 2: C$\to$B$\to$A}{}\\\hline
\rule{0pt}{8mm}\textbf{III}\ \begin{tikzpicture}[baseline=(k3.base),x=0.9cm,y=0.6cm]
\foreach \n/\x/\y in {1/0/1,2/2/1,3/1/0,4/3/0} {\node[kknoten,minimum size=5mm,inner sep=0pt,font=\tiny] (k\n) at (\x,\y) {K\n};}
\draw[kante] (k1) -- node[gewicht] {4} (k2); \draw[kante] (k2) -- node[gewicht] {1} (k3);
\draw[kante] (k3) -- node[gewicht] {5} (k1); \draw[kante] (k3) -- node[gewicht] {2} (k4);
\end{tikzpicture}
& \lsgoder{\lsgfarbe stark}{} & \lsgoder{\lsgfarbe nein}{} & \lsgoder{\lsgfarbe ja}{} & \lsgoder{\lsgfarbe ja}{} & \lsgoder{\lsgfarbe 11: K2--K1--K3--K4}{}\\\hline
\end{tabularx}
\end{aufgabe}
\end{document}
"
%           (füllt alle Lücken rot aus; siehe \lsg, \lsglinien, \lsgoder)
% =====================================================================
\usepackage[a4paper,left=18mm,right=18mm,top=26mm,bottom=17mm,headheight=15mm,headsep=4mm,footskip=8mm]{geometry}
\usepackage{iftex}
\ifPDFTeX
  \usepackage[T1]{fontenc}
  \usepackage[utf8]{inputenc}
  \usepackage{helvet}
  \usepackage[scaled=0.86]{beramono}
  \renewcommand{\familydefault}{\sfdefault}
\else
  \usepackage{fontspec}
  \setmainfont{texgyreheros}[Extension=.otf,UprightFont=*-regular,BoldFont=*-bold,ItalicFont=*-italic,BoldItalicFont=*-bolditalic]
  \setmonofont{DejaVuSansMono}[Extension=.ttf,UprightFont=*,BoldFont=*-Bold,ItalicFont=*-Oblique,BoldItalicFont=*-BoldOblique,Scale=0.86]
\fi
\usepackage[ngerman,shorthands=off]{babel}
\usepackage{xcolor,graphicx,tabularx,array,enumitem,fancyhdr,tikz,listings,multicol,calc,etoolbox}
\usepackage[most]{tcolorbox}
\usetikzlibrary{arrows.meta,positioning,calc,decorations.pathreplacing,shapes.geometric}
\usepackage[hidelinks]{hyperref}
\setlength{\parindent}{0pt}
\setlength{\parskip}{3pt}
\setlist{nosep,leftmargin=*}

% ---------- Lösungsschalter ----------
\newif\ifloesung
\ifdefined\mitloesung\loesungtrue\fi
\newcommand{\lsgfarbe}{\color{red!75!black}}

% ---------- Kopf- und Fußzeile (einheitliche Form) ----------
\newcommand{\lehrkraft}{Hau}
\newcommand{\themenbereich}{3. Graphen}
\newcommand{\abnummer}{}
\pagestyle{fancy}\fancyhf{}
\renewcommand{\headrulewidth}{0.4pt}
\fancyhead[L]{\small Klasse 11\\Datum:}
\fancyhead[C]{\small Themenbereich:\\\themenbereich}
\fancyhead[R]{\small \lehrkraft\ifloesung\\{\lsgfarbe\bfseries Lösung}\fi}
\fancyfoot[C]{\scriptsize edu-mrh.de\,·\,Informatik 11\,·\,Blatt \abnummer\ifloesung\ (Lösung)\fi\,·\,Seite \thepage}
\newcommand{\abtitel}[2]{\renewcommand{\abnummer}{#1}\begin{center}\Large\bfseries #1\quad\underline{#2}\end{center}\vspace{-1mm}}

% ---------- Boxen ----------
\newenvironment{aufgabe}[2][]{\begin{tcolorbox}[enhanced,breakable,colback=white,colframe=black,boxrule=0.7pt,arc=3mm,left=2mm,right=2mm,top=1.5mm,bottom=1.5mm,before skip=2.5mm,after skip=2.5mm]\textbf{Aufgabe #2)}\ifstrempty{#1}{}{ \textit{#1}}\par\vspace{0.6mm}}{\end{tcolorbox}}
\newtcolorbox{merke}[1][Merke]{enhanced,breakable,colback=green!5,colframe=green!40!black,boxrule=0.8pt,arc=2mm,left=2mm,right=2mm,top=1.5mm,bottom=1.5mm,before skip=2.5mm,after skip=2.5mm,title={#1},fonttitle=\bfseries,coltitle=white}
\newtcolorbox{hinweis}{enhanced,breakable,colback=yellow!12,colframe=orange!70!black,boxrule=0.6pt,arc=2mm,left=2mm,right=2mm,top=1mm,bottom=1mm,before skip=2mm,after skip=2mm}
\newtcolorbox{ausblick}[1][Ausblick]{enhanced,breakable,colback=teal!5,colframe=teal!60!black,boxrule=0.8pt,arc=2mm,left=2mm,right=2mm,top=1.5mm,bottom=1.5mm,before skip=2.5mm,after skip=2.5mm,title={#1},fonttitle=\bfseries,coltitle=white}
\newcommand{\web}[1]{{\small\color{gray}\textit{Website: #1}}}

% ---------- Lücken, Linien, Karos ----------
\newcommand{\luecke}[1][3.5cm]{\rule[-0.5ex]{#1}{0.4pt}}
\newcommand{\linien}[1]{\par\vspace{1.5mm}\foreach \i in {1,...,#1}{\noindent\rule[0pt]{\linewidth}{0.3pt}\par\vspace{5.5mm}}\vspace{-3mm}}
\newcommand{\kariert}[1]{\par\vspace{1.5mm}\noindent\begin{tikzpicture}\draw[step=5mm,gray!55,thin] (0,0) grid (\linewidth-1pt,#1*5mm);\end{tikzpicture}\par}
\newcommand{\karobox}[2]{\begin{tikzpicture}\draw[step=5mm,gray!55,thin] (0,0) grid (#1,#2);\end{tikzpicture}}
\newcommand{\code}[1]{\texttt{#1}}

% Lücke mit Lösung: \lsg[Mindestbreite]{Lösungstext}
\newlength{\lsgbreite}
\newcommand{\lsg}[2][3.5cm]{\ifloesung\settowidth{\lsgbreite}{\,#2\,}\ifdim\lsgbreite<#1\setlength{\lsgbreite}{#1}\fi\underline{\makebox[\lsgbreite][l]{\lsgfarbe\,#2}}\else\luecke[#1]\fi}
% Schreiblinien mit Lösung: \lsglinien{Anzahl}{Lösungstext} (gleiche Höhe wie \linien)
\newcommand{\lsglinien}[2]{\ifloesung\par\vspace{1.5mm}\noindent\begin{minipage}[t][#1\dimexpr 6.9mm\relax][t]{\linewidth}\lsgfarbe #2\end{minipage}\par\vspace{-1.5mm}\else\linien{#1}\fi}
% Beliebiger Inhalt: \lsgoder{nur in der Lösung}{nur im Blatt}
\newcommand{\lsgoder}[2]{\ifloesung#1\else#2\fi}

% ---------- Graphen zeichnen ----------
\tikzset{
  knoten/.style={draw,rounded corners=3pt,fill=orange!12,minimum height=6mm,inner sep=3pt,font=\small},
  kknoten/.style={draw,circle,fill=orange!12,minimum size=7mm,inner sep=1pt,font=\small},
  kante/.style={thick},
  gkante/.style={thick,-{Stealth[length=2.6mm]}},
  dkante/.style={thick,{Stealth[length=2.6mm]}-{Stealth[length=2.6mm]}},
  gewicht/.style={font=\scriptsize,fill=white,inner sep=1pt,text=teal!60!black},
  lsgkante/.style={thick,red!75!black},
  lsggkante/.style={thick,red!75!black,-{Stealth[length=2.6mm]}},
  lsgknoten/.style={knoten,draw=red!75!black,text=red!75!black,fill=red!4},
}
% Adjazenzmatrix zum Ausfüllen: \adjmx[Lösungseinträge]{Anzahl}{Knotenliste}{Zellbreite}
% Lösungseinträge als Zeile/Spalte/Text, z. B. \adjmx[1/2/X,2/1/X]{3}{A,B,C}{7mm}
\newcommand{\adjmx}[4][]{\begin{tikzpicture}[x=#4,y=#4]
  \foreach \n [count=\i] in {#3} {\node[font=\small\bfseries] at (\i,0) {\n}; \node[font=\small\bfseries,anchor=east] at (0.45,-\i) {\n};}
  \foreach \k in {0,...,#2} {\draw[gray] (\k+0.5,-0.5) -- (\k+0.5,-#2-0.5); \draw[gray] (0.5,-\k-0.5) -- (#2+0.5,-\k-0.5);}
  \ifloesung\ifstrempty{#1}{}{\foreach \z/\s/\t in {#1} {\node[text=red!75!black,font=\small] at (\s,-\z) {\t};}}\fi
\end{tikzpicture}}

% ---------- Java-Listings ----------
\lstset{language=Java,basicstyle=\ttfamily\small,keywordstyle=\bfseries,commentstyle=\color{gray!80!black},stringstyle=\color{black},showstringspaces=false,columns=flexible,keepspaces=true,tabsize=3,frame=single,framerule=0.4pt,rulecolor=\color{black},xleftmargin=2mm,framexleftmargin=1mm,aboveskip=2mm,belowskip=2mm,breaklines=true,escapechar=§,
 literate={ä}{{\"a}}1 {ö}{{\"o}}1 {ü}{{\"u}}1 {Ä}{{\"A}}1 {Ö}{{\"O}}1 {Ü}{{\"U}}1 {ß}{{\ss}}1 {–}{{--}}1 {„}{{\glqq}}1 {“}{{\grqq}}1}
\lstnewenvironment{java}[1][]{\lstset{#1}}{}

%  Übersetzen mit XeLaTeX, LuaLaTeX, pdfLaTeX oder Tectonic.
%
%  Blatt:   xelatex ab-3-1-grundbegriffe.tex
%  Lösung:  xelatex -jobname=ab-3-1-grundbegriffe-lsg "\def\mitloesung{}\documentclass[11pt,a4paper]{article}
% =====================================================================
%  Gemeinsame Vorlage der Arbeitsblätter "Graphen" – Informatik 11 (NTG)
%  (gleiches Layout wie informatik/10/java/arbeitsblaetter/ab-vorlage.tex)
%  Einbinden in jedem Blatt direkt nach \documentclass: \input{ab-vorlage}
%  Übersetzen mit XeLaTeX, LuaLaTeX, pdfLaTeX oder Tectonic.
%
%  Blatt:   xelatex ab-3-1-grundbegriffe.tex
%  Lösung:  xelatex -jobname=ab-3-1-grundbegriffe-lsg "\def\mitloesung{}\input{ab-3-1-grundbegriffe}"
%           (füllt alle Lücken rot aus; siehe \lsg, \lsglinien, \lsgoder)
% =====================================================================
\usepackage[a4paper,left=18mm,right=18mm,top=26mm,bottom=17mm,headheight=15mm,headsep=4mm,footskip=8mm]{geometry}
\usepackage{iftex}
\ifPDFTeX
  \usepackage[T1]{fontenc}
  \usepackage[utf8]{inputenc}
  \usepackage{helvet}
  \usepackage[scaled=0.86]{beramono}
  \renewcommand{\familydefault}{\sfdefault}
\else
  \usepackage{fontspec}
  \setmainfont{texgyreheros}[Extension=.otf,UprightFont=*-regular,BoldFont=*-bold,ItalicFont=*-italic,BoldItalicFont=*-bolditalic]
  \setmonofont{DejaVuSansMono}[Extension=.ttf,UprightFont=*,BoldFont=*-Bold,ItalicFont=*-Oblique,BoldItalicFont=*-BoldOblique,Scale=0.86]
\fi
\usepackage[ngerman,shorthands=off]{babel}
\usepackage{xcolor,graphicx,tabularx,array,enumitem,fancyhdr,tikz,listings,multicol,calc,etoolbox}
\usepackage[most]{tcolorbox}
\usetikzlibrary{arrows.meta,positioning,calc,decorations.pathreplacing,shapes.geometric}
\usepackage[hidelinks]{hyperref}
\setlength{\parindent}{0pt}
\setlength{\parskip}{3pt}
\setlist{nosep,leftmargin=*}

% ---------- Lösungsschalter ----------
\newif\ifloesung
\ifdefined\mitloesung\loesungtrue\fi
\newcommand{\lsgfarbe}{\color{red!75!black}}

% ---------- Kopf- und Fußzeile (einheitliche Form) ----------
\newcommand{\lehrkraft}{Hau}
\newcommand{\themenbereich}{3. Graphen}
\newcommand{\abnummer}{}
\pagestyle{fancy}\fancyhf{}
\renewcommand{\headrulewidth}{0.4pt}
\fancyhead[L]{\small Klasse 11\\Datum:}
\fancyhead[C]{\small Themenbereich:\\\themenbereich}
\fancyhead[R]{\small \lehrkraft\ifloesung\\{\lsgfarbe\bfseries Lösung}\fi}
\fancyfoot[C]{\scriptsize edu-mrh.de\,·\,Informatik 11\,·\,Blatt \abnummer\ifloesung\ (Lösung)\fi\,·\,Seite \thepage}
\newcommand{\abtitel}[2]{\renewcommand{\abnummer}{#1}\begin{center}\Large\bfseries #1\quad\underline{#2}\end{center}\vspace{-1mm}}

% ---------- Boxen ----------
\newenvironment{aufgabe}[2][]{\begin{tcolorbox}[enhanced,breakable,colback=white,colframe=black,boxrule=0.7pt,arc=3mm,left=2mm,right=2mm,top=1.5mm,bottom=1.5mm,before skip=2.5mm,after skip=2.5mm]\textbf{Aufgabe #2)}\ifstrempty{#1}{}{ \textit{#1}}\par\vspace{0.6mm}}{\end{tcolorbox}}
\newtcolorbox{merke}[1][Merke]{enhanced,breakable,colback=green!5,colframe=green!40!black,boxrule=0.8pt,arc=2mm,left=2mm,right=2mm,top=1.5mm,bottom=1.5mm,before skip=2.5mm,after skip=2.5mm,title={#1},fonttitle=\bfseries,coltitle=white}
\newtcolorbox{hinweis}{enhanced,breakable,colback=yellow!12,colframe=orange!70!black,boxrule=0.6pt,arc=2mm,left=2mm,right=2mm,top=1mm,bottom=1mm,before skip=2mm,after skip=2mm}
\newtcolorbox{ausblick}[1][Ausblick]{enhanced,breakable,colback=teal!5,colframe=teal!60!black,boxrule=0.8pt,arc=2mm,left=2mm,right=2mm,top=1.5mm,bottom=1.5mm,before skip=2.5mm,after skip=2.5mm,title={#1},fonttitle=\bfseries,coltitle=white}
\newcommand{\web}[1]{{\small\color{gray}\textit{Website: #1}}}

% ---------- Lücken, Linien, Karos ----------
\newcommand{\luecke}[1][3.5cm]{\rule[-0.5ex]{#1}{0.4pt}}
\newcommand{\linien}[1]{\par\vspace{1.5mm}\foreach \i in {1,...,#1}{\noindent\rule[0pt]{\linewidth}{0.3pt}\par\vspace{5.5mm}}\vspace{-3mm}}
\newcommand{\kariert}[1]{\par\vspace{1.5mm}\noindent\begin{tikzpicture}\draw[step=5mm,gray!55,thin] (0,0) grid (\linewidth-1pt,#1*5mm);\end{tikzpicture}\par}
\newcommand{\karobox}[2]{\begin{tikzpicture}\draw[step=5mm,gray!55,thin] (0,0) grid (#1,#2);\end{tikzpicture}}
\newcommand{\code}[1]{\texttt{#1}}

% Lücke mit Lösung: \lsg[Mindestbreite]{Lösungstext}
\newlength{\lsgbreite}
\newcommand{\lsg}[2][3.5cm]{\ifloesung\settowidth{\lsgbreite}{\,#2\,}\ifdim\lsgbreite<#1\setlength{\lsgbreite}{#1}\fi\underline{\makebox[\lsgbreite][l]{\lsgfarbe\,#2}}\else\luecke[#1]\fi}
% Schreiblinien mit Lösung: \lsglinien{Anzahl}{Lösungstext} (gleiche Höhe wie \linien)
\newcommand{\lsglinien}[2]{\ifloesung\par\vspace{1.5mm}\noindent\begin{minipage}[t][#1\dimexpr 6.9mm\relax][t]{\linewidth}\lsgfarbe #2\end{minipage}\par\vspace{-1.5mm}\else\linien{#1}\fi}
% Beliebiger Inhalt: \lsgoder{nur in der Lösung}{nur im Blatt}
\newcommand{\lsgoder}[2]{\ifloesung#1\else#2\fi}

% ---------- Graphen zeichnen ----------
\tikzset{
  knoten/.style={draw,rounded corners=3pt,fill=orange!12,minimum height=6mm,inner sep=3pt,font=\small},
  kknoten/.style={draw,circle,fill=orange!12,minimum size=7mm,inner sep=1pt,font=\small},
  kante/.style={thick},
  gkante/.style={thick,-{Stealth[length=2.6mm]}},
  dkante/.style={thick,{Stealth[length=2.6mm]}-{Stealth[length=2.6mm]}},
  gewicht/.style={font=\scriptsize,fill=white,inner sep=1pt,text=teal!60!black},
  lsgkante/.style={thick,red!75!black},
  lsggkante/.style={thick,red!75!black,-{Stealth[length=2.6mm]}},
  lsgknoten/.style={knoten,draw=red!75!black,text=red!75!black,fill=red!4},
}
% Adjazenzmatrix zum Ausfüllen: \adjmx[Lösungseinträge]{Anzahl}{Knotenliste}{Zellbreite}
% Lösungseinträge als Zeile/Spalte/Text, z. B. \adjmx[1/2/X,2/1/X]{3}{A,B,C}{7mm}
\newcommand{\adjmx}[4][]{\begin{tikzpicture}[x=#4,y=#4]
  \foreach \n [count=\i] in {#3} {\node[font=\small\bfseries] at (\i,0) {\n}; \node[font=\small\bfseries,anchor=east] at (0.45,-\i) {\n};}
  \foreach \k in {0,...,#2} {\draw[gray] (\k+0.5,-0.5) -- (\k+0.5,-#2-0.5); \draw[gray] (0.5,-\k-0.5) -- (#2+0.5,-\k-0.5);}
  \ifloesung\ifstrempty{#1}{}{\foreach \z/\s/\t in {#1} {\node[text=red!75!black,font=\small] at (\s,-\z) {\t};}}\fi
\end{tikzpicture}}

% ---------- Java-Listings ----------
\lstset{language=Java,basicstyle=\ttfamily\small,keywordstyle=\bfseries,commentstyle=\color{gray!80!black},stringstyle=\color{black},showstringspaces=false,columns=flexible,keepspaces=true,tabsize=3,frame=single,framerule=0.4pt,rulecolor=\color{black},xleftmargin=2mm,framexleftmargin=1mm,aboveskip=2mm,belowskip=2mm,breaklines=true,escapechar=§,
 literate={ä}{{\"a}}1 {ö}{{\"o}}1 {ü}{{\"u}}1 {Ä}{{\"A}}1 {Ö}{{\"O}}1 {Ü}{{\"U}}1 {ß}{{\ss}}1 {–}{{--}}1 {„}{{\glqq}}1 {“}{{\grqq}}1}
\lstnewenvironment{java}[1][]{\lstset{#1}}{}

% Dörfer-Graph (Merkkasten Begriffe): Koordinaten wie im Begriffe-Explorer der Website
\newcommand{\doerferknoten}{%
  \node[knoten,font=\tiny,minimum height=4mm,inner sep=1.5pt] (A) at (4.4,3.3) {Altdorf};
  \node[knoten,font=\tiny,minimum height=4mm,inner sep=1.5pt] (W) at (1.2,2.9) {Weiler};
  \node[knoten,font=\tiny,minimum height=4mm,inner sep=1.5pt] (F) at (5.1,2.0) {Fischbach};
  \node[knoten,font=\tiny,minimum height=4mm,inner sep=1.5pt] (Z) at (2.9,2.0) {Ziegelstein};
  \node[knoten,font=\tiny,minimum height=4mm,inner sep=1.5pt] (N) at (0.9,0.95) {Neustadt};
  \node[knoten,font=\tiny,minimum height=4mm,inner sep=1.5pt] (B) at (4.4,0.8) {Burg};
  \node[knoten,font=\tiny,minimum height=4mm,inner sep=1.5pt] (R) at (2.6,0.2) {Rain};}
\begin{document}
\abtitel{3.1}{Graphen -- Grundbegriffe}

Hänsel und Gretel haben die Hexe überlistet. Doch die Brotkrumen, mit denen sie ihren Weg markiert hatten, haben die Vögel gefressen. Findest du trotzdem nach Hause?

\begin{aufgabe}[Hänsel und Gretel]{1}
Starte das Spiel auf der Website (Station 3.1).
\textbf{a)} Zeichne beim Spielen eine Karte des Waldes: jeder Ort ein Kästchen, jeder Weg eine Linie. Führt ein Weg nur in eine Richtung, zeichne einen Pfeil. Das Hexenhaus ist schon eingezeichnet.
\end{aufgabe}
\noindent\begin{tikzpicture}
\draw[step=5mm,gray!55,thin] (0,0) grid (17.4,7);
\node[knoten,fill=green!15] (hex) at (4,2.3) {Hexenhaus (Start)};
\ifloesung
  \node[lsgknoten] (warn) at (5.5,6.5) {Warnung};
  \node[lsgknoten] (wolf) at (10.5,6.5) {Werwolf};
  \node[lsgknoten] (stadt) at (1.6,5.2) {Stadt};
  \node[lsgknoten] (abz) at (5.5,5.2) {Abzweigung};
  \node[lsgknoten] (spinne) at (13,5.2) {Spinne};
  \node[lsgknoten] (baeck) at (1.6,3.6) {Bäckerei};
  \node[lsgknoten] (pilze) at (4,3.9) {Pilze};
  \node[lsgknoten] (w3) at (13,3.4) {Wald3};
  \node[lsgknoten] (w1) at (8,1.0) {Wald1};
  \node[lsgknoten] (stein) at (12.2,2.0) {Steinkreis};
  \node[lsgknoten] (w2) at (15,0.8) {Wald2};
  \node[lsgknoten] (heim) at (16.2,2.6) {Zuhause};
  \draw[lsgkante] (stadt) -- (baeck);
  \draw[lsgkante] (stadt) -- (pilze);
  \draw[lsgkante] (pilze) -- (abz);
  \draw[lsgkante] (pilze) -- (hex);
  \draw[lsggkante] (pilze) to[bend left=50] node[right,font=\scriptsize,text=red!75!black] {Pilze essen} (hex);
  \draw[lsgkante] (abz) -- (warn);
  \draw[lsggkante] (warn) -- (wolf);
  \draw[lsgkante] (abz) -- (spinne);
  \draw[lsgkante] (spinne) -- (w3);
  \draw[lsgkante] (hex) -- (w1);
  \draw[lsgkante] (w1) -- (w3);
  \draw[lsgkante] (w1) -- (w2);
  \draw[lsgkante] (w1) -- (stein);
  \draw[lsgkante] (w2) -- (w3);
  \draw[lsgkante] (w2) -- (stein);
  \draw[lsgkante] (w3) -- (stein);
  \draw[lsggkante] (w2) -- (heim);
  \node[text=red!75!black,font=\scriptsize,align=left,anchor=north west] at (0.2,1.3) {13 Knoten, 17 Kanten\\(14 ungerichtet, 3 gerichtet)};
\fi
\end{tikzpicture}

\textbf{b)} Am Ende zeigt dir das Spiel deinen Weg. Notiere ihn und die Anzahl der Schritte:
\lsglinien{1}{individuell; kürzester Weg: Hexenhaus -- Wald1 -- Wald2 -- Zuhause (3 Schritte)}

\begin{minipage}[t]{0.6\linewidth}
\begin{merke}[Graph und Kurzschreibweise $G = (V, E)$]\raggedright
Ein Graph besteht aus \lsg[2.3cm]{Knoten} (engl. \textit{vertices}) und \lsg[2.3cm]{Kanten} (engl. \textit{edges}). Eine Kante verbindet zwei Knoten. Im Spiel sind die Orte die \lsg[2cm]{Knoten}, die Wege die \lsg[2cm]{Kanten}.\\[1mm]
Kurzschreibweise $G = (V, E)$: $V$ ist die Menge der Knoten, $E$ die Menge der Kanten.
\begin{itemize}
\item $(A, B)$: \lsg[2.5cm]{gerichtete} Kante, nur von $A$ nach $B$
\item $\{A, B\}$: \lsg[2.5cm]{ungerichtete} Kante, steht für $(A, B)$ und \lsg[1.4cm]{$(B, A)$}
\end{itemize}
\end{merke}
\end{minipage}\hfill
\begin{minipage}[t]{0.37\linewidth}
\vspace{1mm}
\centering\small \textbf{Beispiel: Wer folgt wem?}\\[1mm]
\begin{tikzpicture}[x=0.95cm,y=0.9cm]
\node[knoten] (leon) at (0.7,2.2) {Leon};
\node[knoten] (letifa) at (3.3,2.4) {Letifa};
\node[knoten] (luise) at (4.3,1.1) {Luise};
\node[knoten] (hamza) at (2.4,0.3) {Hamza};
\node[knoten] (sabine) at (0.5,0.6) {Sabine};
\draw[dkante] (leon) -- (letifa);
\draw[gkante] (letifa) -- (luise);
\draw[gkante] (leon) -- (luise);
\draw[dkante] (leon) -- (hamza);
\draw[gkante] (hamza) -- (luise);
\end{tikzpicture}
\end{minipage}\\[1mm]
{\small $V = \{\text{Hamza}, \text{Leon}, \text{Letifa}, \text{Luise}, \text{Sabine}\}$\\
$E = \{\{\text{Leon}, \text{Letifa}\}, \{\text{Leon}, \text{Hamza}\}, (\text{Letifa}, \text{Luise}), (\text{Leon}, \text{Luise}), (\text{Hamza}, \text{Luise})\}$}

\begin{aufgabe}[Hänsel und Gretel]{1}
\textbf{c)} Schreibe den Ausschnitt deiner Karte mit den Knoten Hexenhaus, Pilze, Stadt, Bäckerei und Wald1 in Kurzschreibweise auf (5 Kanten).\\[2mm]
$V = \{$\lsg[15.4cm]{Hexenhaus, Pilze, Stadt, Bäckerei, Wald1}$\}$\\[3mm]
$E = \{$\lsg[15.4cm]{\{Hexenhaus, Pilze\}, (Pilze, Hexenhaus), \{Pilze, Stadt\}, \{Stadt, Bäckerei\}, \{Hexenhaus, Wald1\}}$\}$
\end{aufgabe}

\newpage
\begin{minipage}[t]{0.71\linewidth}
\begin{merke}[Pfade, Zyklen, Richtung, Gewicht]\raggedright\small
\textbf{Pfad (Weg):} Folge von Knoten, in der aufeinanderfolgende Knoten durch eine \lsg[1.8cm]{Kante} verbunden sind; keine Kante wird doppelt benutzt.\\[0.6mm]
\textbf{Einfacher Pfad:} Pfad, in dem sich kein \lsg[1.8cm]{Knoten} wiederholt.\\
Beispiel in A: \lsg[5.6cm]{Altdorf -- Weiler -- Ziegelstein -- Neustadt}\\[0.6mm]
\textbf{Länge eines Pfades:} ungewichtet die \lsg[3.2cm]{Anzahl der Kanten}, gewichtet die \lsg[3.2cm]{Summe der Gewichte}.\\[0.6mm]
\textbf{Zyklus:} einfacher Pfad, dessen \lsg[5cm]{Start- und Endknoten gleich sind}.\\
Beispiel in A: \lsg[5.6cm]{Weiler -- Ziegelstein -- Fischbach -- Weiler}\\
Graph mit Zyklus: \emph{zyklisch}, sonst \lsg[3.4cm]{azyklisch (zyklenfrei)}.\\[0.6mm]
\textbf{Gerichtet:} Kanten nur \lsg[3cm]{in Pfeilrichtung} (Pfeil).\ \textbf{Ungerichtet:} \lsg[2.8cm]{in beide Richtungen} (Strich).\\[0.6mm]
\textbf{Stark zusammenhängend:} Von \lsg[2.4cm]{jedem Knoten} gibt es einen Pfad zu \lsg[3.2cm]{jedem anderen Knoten}.\\
\textbf{Schwach zusammenhängend:} \lsg[5.4cm]{nur ohne Beachtung der Richtungen}\\ \lsg[\linewidth]{hängt alles zusammen. B ist schwach, denn von Rain führt keine Kante weg.}\\[0.6mm]
\textbf{Gewichtet:} Jede Kante trägt eine \lsg[2.6cm]{Zahl (Gewicht)}, z.\,B. \lsg[3.6cm]{Entfernung, Fahrzeit}.
\end{merke}
\end{minipage}\hfill
\begin{minipage}[t]{0.27\linewidth}
\vspace{1mm}
\textbf{\small A} {\scriptsize ungerichtet}\\
\begin{tikzpicture}[x=0.8cm,y=0.85cm]
\doerferknoten
\draw[kante] (A)--(W); \draw[kante] (F)--(A); \draw[kante] (W)--(F); \draw[kante] (W)--(Z);
\draw[kante] (Z)--(F); \draw[kante] (Z)--(N); \draw[kante] (B)--(Z); \draw[kante] (F)--(B);
\draw[kante] (N)--(R); \draw[kante] (B)--(R);
\end{tikzpicture}\\[2mm]
\textbf{\small B} {\scriptsize gerichtet}\\
\begin{tikzpicture}[x=0.8cm,y=0.85cm]
\doerferknoten
\draw[gkante] (A)--(W); \draw[gkante] (F)--(A); \draw[gkante] (W)--(F); \draw[gkante] (W)--(Z);
\draw[gkante] (Z)--(F); \draw[gkante] (Z)--(N); \draw[gkante] (B)--(Z); \draw[gkante] (F)--(B);
\draw[gkante] (N)--(R); \draw[gkante] (B)--(R);
\end{tikzpicture}
\end{minipage}

\begin{aufgabe}{2}
{\raggedright Zeichne $G = (V, E)$ mit $V = \{A, B, C, D, E, F\}$ und\\ $E = \{(A, B), (A, C), (A, E), (B, C), (C, C), (C, D), (E, C), (F, D)\}$. Beantworte die Fragen.\par}\vspace{1mm}
\begin{minipage}[t]{0.42\linewidth}
\vspace{0pt}
\begin{tikzpicture}
\draw[step=5mm,gray!55,thin] (0,0) grid (7,3);
\ifloesung
  \foreach \n/\x/\y/\t in {a/0.6/1.5/A, b/2/2.55/B, e/2/0.45/E, c/3.4/1.5/C, d/4.9/1.5/D, f/6.4/1.5/F} {
    \node[kknoten,minimum size=6mm,draw=red!75!black,text=red!75!black] (\n) at (\x,\y) {\t};}
  \draw[lsggkante] (a)--(b); \draw[lsggkante] (a)--(c); \draw[lsggkante] (a)--(e);
  \draw[lsggkante] (b)--(c); \draw[lsggkante] (c)--(d); \draw[lsggkante] (e)--(c);
  \draw[lsggkante] (f)--(d);
  \draw[lsggkante] (c) to[out=60,in=120,looseness=6] (c);
\fi
\end{tikzpicture}
\end{minipage}\hfill
\begin{minipage}[t]{0.55\linewidth}
\vspace{0pt}\small\raggedright
Gerichtet? Woran erkennbar? \lsg[3.9cm]{ja, Paare ( ) statt \{ \}}\\[1.5mm]
Ein Zyklus: \lsg[5.7cm]{(C, C), eine Schlinge (Länge 1)}\\[1.5mm]
Stark/schwach zusammenhängend? \lsg[3.6cm]{schwach: von D}\\[1.5mm]
\lsg[0.98\linewidth]{kommt man nirgends hin.}\\[1.5mm]
Länge des längsten einfachen Pfades: \lsg[3.2cm]{3 (A--B--C--D)}
\end{minipage}
\end{aufgabe}

\begin{aufgabe}[Zurück in den Wald]{3}
Deine Karte aus Aufgabe 1: \textbf{a)} gerichtet?\ \textbf{b)} ein Zyklus?\ \textbf{c)} stark oder schwach zusammenhängend?\ \textbf{d)} „Der Weg zurück … scheint verschwunden“: Ist die Kante Hexenhaus -- Wald1 gerichtet? Prüfe im Spiel.
\lsglinien{2}{\small a) gemischt; gerichtet nur (Pilze, Hexenhaus), (Warnung, Werwolf), (Wald2, Zuhause)\quad b) Wald1 -- Wald2 -- Wald3 -- Wald1\\ c) schwach: Von Zuhause/Werwolf führt keine Kante weg.\quad d) nein: Von Wald2 aus führt ein Weg zurück.}
\end{aufgabe}

\begin{aufgabe}[Ordne die Graphen ein.]{4}
\vspace{-1mm}
\small\renewcommand{\arraystretch}{1.15}
\begin{tabularx}{\linewidth}{|>{\centering\arraybackslash}m{5.6cm}|>{\centering\arraybackslash}m{1.7cm}|>{\centering\arraybackslash}m{1.1cm}|>{\centering\arraybackslash}m{1.1cm}|>{\centering\arraybackslash}m{1.1cm}|>{\centering\arraybackslash}X|}\hline
\textbf{Graph} & \scriptsize\textbf{zusammen-\newline hängend?}\newline stark/schwach/nein & \scriptsize\textbf{ge-\newline richtet?} & \scriptsize\textbf{ge-\newline wichtet?} & \scriptsize\textbf{zyklisch?} & \scriptsize\textbf{Länge eines längsten einfachen Pfades}\\\hline
\rule{0pt}{8mm}\textbf{I}\ \begin{tikzpicture}[baseline=(k5.base),x=0.85cm,y=0.55cm]
\foreach \n/\x/\y in {1/0/1,2/1/1,3/2/1,4/3/1,5/1/0,6/3/0} {\node[kknoten,minimum size=5mm,inner sep=0pt,font=\tiny] (k\n) at (\x,\y) {K\n};}
\draw[kante] (k1)--(k2)--(k3)--(k4); \draw[kante] (k2)--(k5);
\end{tikzpicture}
& \lsgoder{\lsgfarbe nein}{} & \lsgoder{\lsgfarbe nein}{} & \lsgoder{\lsgfarbe nein}{} & \lsgoder{\lsgfarbe nein}{} & \lsgoder{\lsgfarbe 3: K1--K2--K3--K4}{}\\\hline
\rule{0pt}{5mm}\textbf{II}\ $V = \{A, B, C\}$\newline $E =\{(A, B), (B, A), (C, B), (C, C)\}$
& \lsgoder{\lsgfarbe schwach}{} & \lsgoder{\lsgfarbe ja}{} & \lsgoder{\lsgfarbe nein}{} & \lsgoder{\lsgfarbe ja}{} & \lsgoder{\lsgfarbe 2: C$\to$B$\to$A}{}\\\hline
\rule{0pt}{8mm}\textbf{III}\ \begin{tikzpicture}[baseline=(k3.base),x=0.9cm,y=0.6cm]
\foreach \n/\x/\y in {1/0/1,2/2/1,3/1/0,4/3/0} {\node[kknoten,minimum size=5mm,inner sep=0pt,font=\tiny] (k\n) at (\x,\y) {K\n};}
\draw[kante] (k1) -- node[gewicht] {4} (k2); \draw[kante] (k2) -- node[gewicht] {1} (k3);
\draw[kante] (k3) -- node[gewicht] {5} (k1); \draw[kante] (k3) -- node[gewicht] {2} (k4);
\end{tikzpicture}
& \lsgoder{\lsgfarbe stark}{} & \lsgoder{\lsgfarbe nein}{} & \lsgoder{\lsgfarbe ja}{} & \lsgoder{\lsgfarbe ja}{} & \lsgoder{\lsgfarbe 11: K2--K1--K3--K4}{}\\\hline
\end{tabularx}
\end{aufgabe}
\end{document}
"
%           (füllt alle Lücken rot aus; siehe \lsg, \lsglinien, \lsgoder)
% =====================================================================
\usepackage[a4paper,left=18mm,right=18mm,top=26mm,bottom=17mm,headheight=15mm,headsep=4mm,footskip=8mm]{geometry}
\usepackage{iftex}
\ifPDFTeX
  \usepackage[T1]{fontenc}
  \usepackage[utf8]{inputenc}
  \usepackage{helvet}
  \usepackage[scaled=0.86]{beramono}
  \renewcommand{\familydefault}{\sfdefault}
\else
  \usepackage{fontspec}
  \setmainfont{texgyreheros}[Extension=.otf,UprightFont=*-regular,BoldFont=*-bold,ItalicFont=*-italic,BoldItalicFont=*-bolditalic]
  \setmonofont{DejaVuSansMono}[Extension=.ttf,UprightFont=*,BoldFont=*-Bold,ItalicFont=*-Oblique,BoldItalicFont=*-BoldOblique,Scale=0.86]
\fi
\usepackage[ngerman,shorthands=off]{babel}
\usepackage{xcolor,graphicx,tabularx,array,enumitem,fancyhdr,tikz,listings,multicol,calc,etoolbox}
\usepackage[most]{tcolorbox}
\usetikzlibrary{arrows.meta,positioning,calc,decorations.pathreplacing,shapes.geometric}
\usepackage[hidelinks]{hyperref}
\setlength{\parindent}{0pt}
\setlength{\parskip}{3pt}
\setlist{nosep,leftmargin=*}

% ---------- Lösungsschalter ----------
\newif\ifloesung
\ifdefined\mitloesung\loesungtrue\fi
\newcommand{\lsgfarbe}{\color{red!75!black}}

% ---------- Kopf- und Fußzeile (einheitliche Form) ----------
\newcommand{\lehrkraft}{Hau}
\newcommand{\themenbereich}{3. Graphen}
\newcommand{\abnummer}{}
\pagestyle{fancy}\fancyhf{}
\renewcommand{\headrulewidth}{0.4pt}
\fancyhead[L]{\small Klasse 11\\Datum:}
\fancyhead[C]{\small Themenbereich:\\\themenbereich}
\fancyhead[R]{\small \lehrkraft\ifloesung\\{\lsgfarbe\bfseries Lösung}\fi}
\fancyfoot[C]{\scriptsize edu-mrh.de\,·\,Informatik 11\,·\,Blatt \abnummer\ifloesung\ (Lösung)\fi\,·\,Seite \thepage}
\newcommand{\abtitel}[2]{\renewcommand{\abnummer}{#1}\begin{center}\Large\bfseries #1\quad\underline{#2}\end{center}\vspace{-1mm}}

% ---------- Boxen ----------
\newenvironment{aufgabe}[2][]{\begin{tcolorbox}[enhanced,breakable,colback=white,colframe=black,boxrule=0.7pt,arc=3mm,left=2mm,right=2mm,top=1.5mm,bottom=1.5mm,before skip=2.5mm,after skip=2.5mm]\textbf{Aufgabe #2)}\ifstrempty{#1}{}{ \textit{#1}}\par\vspace{0.6mm}}{\end{tcolorbox}}
\newtcolorbox{merke}[1][Merke]{enhanced,breakable,colback=green!5,colframe=green!40!black,boxrule=0.8pt,arc=2mm,left=2mm,right=2mm,top=1.5mm,bottom=1.5mm,before skip=2.5mm,after skip=2.5mm,title={#1},fonttitle=\bfseries,coltitle=white}
\newtcolorbox{hinweis}{enhanced,breakable,colback=yellow!12,colframe=orange!70!black,boxrule=0.6pt,arc=2mm,left=2mm,right=2mm,top=1mm,bottom=1mm,before skip=2mm,after skip=2mm}
\newtcolorbox{ausblick}[1][Ausblick]{enhanced,breakable,colback=teal!5,colframe=teal!60!black,boxrule=0.8pt,arc=2mm,left=2mm,right=2mm,top=1.5mm,bottom=1.5mm,before skip=2.5mm,after skip=2.5mm,title={#1},fonttitle=\bfseries,coltitle=white}
\newcommand{\web}[1]{{\small\color{gray}\textit{Website: #1}}}

% ---------- Lücken, Linien, Karos ----------
\newcommand{\luecke}[1][3.5cm]{\rule[-0.5ex]{#1}{0.4pt}}
\newcommand{\linien}[1]{\par\vspace{1.5mm}\foreach \i in {1,...,#1}{\noindent\rule[0pt]{\linewidth}{0.3pt}\par\vspace{5.5mm}}\vspace{-3mm}}
\newcommand{\kariert}[1]{\par\vspace{1.5mm}\noindent\begin{tikzpicture}\draw[step=5mm,gray!55,thin] (0,0) grid (\linewidth-1pt,#1*5mm);\end{tikzpicture}\par}
\newcommand{\karobox}[2]{\begin{tikzpicture}\draw[step=5mm,gray!55,thin] (0,0) grid (#1,#2);\end{tikzpicture}}
\newcommand{\code}[1]{\texttt{#1}}

% Lücke mit Lösung: \lsg[Mindestbreite]{Lösungstext}
\newlength{\lsgbreite}
\newcommand{\lsg}[2][3.5cm]{\ifloesung\settowidth{\lsgbreite}{\,#2\,}\ifdim\lsgbreite<#1\setlength{\lsgbreite}{#1}\fi\underline{\makebox[\lsgbreite][l]{\lsgfarbe\,#2}}\else\luecke[#1]\fi}
% Schreiblinien mit Lösung: \lsglinien{Anzahl}{Lösungstext} (gleiche Höhe wie \linien)
\newcommand{\lsglinien}[2]{\ifloesung\par\vspace{1.5mm}\noindent\begin{minipage}[t][#1\dimexpr 6.9mm\relax][t]{\linewidth}\lsgfarbe #2\end{minipage}\par\vspace{-1.5mm}\else\linien{#1}\fi}
% Beliebiger Inhalt: \lsgoder{nur in der Lösung}{nur im Blatt}
\newcommand{\lsgoder}[2]{\ifloesung#1\else#2\fi}

% ---------- Graphen zeichnen ----------
\tikzset{
  knoten/.style={draw,rounded corners=3pt,fill=orange!12,minimum height=6mm,inner sep=3pt,font=\small},
  kknoten/.style={draw,circle,fill=orange!12,minimum size=7mm,inner sep=1pt,font=\small},
  kante/.style={thick},
  gkante/.style={thick,-{Stealth[length=2.6mm]}},
  dkante/.style={thick,{Stealth[length=2.6mm]}-{Stealth[length=2.6mm]}},
  gewicht/.style={font=\scriptsize,fill=white,inner sep=1pt,text=teal!60!black},
  lsgkante/.style={thick,red!75!black},
  lsggkante/.style={thick,red!75!black,-{Stealth[length=2.6mm]}},
  lsgknoten/.style={knoten,draw=red!75!black,text=red!75!black,fill=red!4},
}
% Adjazenzmatrix zum Ausfüllen: \adjmx[Lösungseinträge]{Anzahl}{Knotenliste}{Zellbreite}
% Lösungseinträge als Zeile/Spalte/Text, z. B. \adjmx[1/2/X,2/1/X]{3}{A,B,C}{7mm}
\newcommand{\adjmx}[4][]{\begin{tikzpicture}[x=#4,y=#4]
  \foreach \n [count=\i] in {#3} {\node[font=\small\bfseries] at (\i,0) {\n}; \node[font=\small\bfseries,anchor=east] at (0.45,-\i) {\n};}
  \foreach \k in {0,...,#2} {\draw[gray] (\k+0.5,-0.5) -- (\k+0.5,-#2-0.5); \draw[gray] (0.5,-\k-0.5) -- (#2+0.5,-\k-0.5);}
  \ifloesung\ifstrempty{#1}{}{\foreach \z/\s/\t in {#1} {\node[text=red!75!black,font=\small] at (\s,-\z) {\t};}}\fi
\end{tikzpicture}}

% ---------- Java-Listings ----------
\lstset{language=Java,basicstyle=\ttfamily\small,keywordstyle=\bfseries,commentstyle=\color{gray!80!black},stringstyle=\color{black},showstringspaces=false,columns=flexible,keepspaces=true,tabsize=3,frame=single,framerule=0.4pt,rulecolor=\color{black},xleftmargin=2mm,framexleftmargin=1mm,aboveskip=2mm,belowskip=2mm,breaklines=true,escapechar=§,
 literate={ä}{{\"a}}1 {ö}{{\"o}}1 {ü}{{\"u}}1 {Ä}{{\"A}}1 {Ö}{{\"O}}1 {Ü}{{\"U}}1 {ß}{{\ss}}1 {–}{{--}}1 {„}{{\glqq}}1 {“}{{\grqq}}1}
\lstnewenvironment{java}[1][]{\lstset{#1}}{}

\begin{document}
\abtitel{3.2}{Graphen implementieren -- die Adjazenzmatrix}

\begin{minipage}[t]{0.6\linewidth}
\vspace{0pt}\raggedright
Beim Ausflug in einen Vergnügungspark werden die Sitzplätze der Achterbahn verteilt. Björn, Doris, Fatma, Gerda, Jakob und Simon haben mit einem Kreuz notiert, mit wem sie gerne fahren möchten. Gelesen wird zeilenweise: Björn (Zeile B) möchte mit Doris und mit Gerda fahren.
\end{minipage}\hfill
\begin{minipage}[t]{0.36\linewidth}
\vspace{0pt}\centering\small
\renewcommand{\arraystretch}{1.1}
\begin{tabular}{|c|c|c|c|c|c|c|}\hline
 & \textbf{B} & \textbf{D} & \textbf{F} & \textbf{G} & \textbf{J} & \textbf{S}\\\hline
\textbf{B} &   & X &   & X &   &  \\\hline
\textbf{D} & X &   &   &   & X & X\\\hline
\textbf{F} &   & X &   & X &   &  \\\hline
\textbf{G} &   &   &   &   & X & X\\\hline
\textbf{J} &   &   & X &   &   &  \\\hline
\textbf{S} & X &   & X &   & X &  \\\hline
\end{tabular}
\end{minipage}

\begin{aufgabe}[Achterbahn]{1}
\begin{minipage}[t]{0.42\linewidth}
\vspace{0pt}
\begin{tikzpicture}
\draw[step=5mm,gray!55,thin] (0,0) grid (7,3);
\ifloesung
  \foreach \n/\x/\y in {B/0.8/2.5, D/3.5/2.55, G/6.2/2.5, S/0.8/0.6, F/3.5/0.5, J/6.2/0.6} {
    \node[kknoten,minimum size=6mm,draw=red!75!black,text=red!75!black] (\n) at (\x,\y) {\n};}
  \draw[lsggkante,{Stealth[length=2.6mm]}-{Stealth[length=2.6mm]}] (B)--(D);
  \draw[lsggkante] (B)--(G); \draw[lsggkante] (D)--(J); \draw[lsggkante] (D)--(S);
  \draw[lsggkante] (F)--(D); \draw[lsggkante] (F)--(G); \draw[lsggkante] (G)--(J);
  \draw[lsggkante] (G)--(S); \draw[lsggkante] (J)--(F); \draw[lsggkante] (S)--(B);
  \draw[lsggkante] (S)--(F); \draw[lsggkante] (S) to[bend right=14] (J);
\fi
\end{tikzpicture}
\end{minipage}\hfill
\begin{minipage}[t]{0.55\linewidth}
\vspace{0pt}\small\raggedright
\textbf{a)} Begründe, dass die Tabelle einen Graphen darstellt, und zeichne ihn links.\\
\lsg[0.98\linewidth]{Personen = Knoten, Kreuz = Kante von Zeile zu Spalte}\\[1.5mm]
\textbf{b)} Ist der Graph gerichtet? Begründe mit einem Beispiel.\\
\lsg[0.98\linewidth]{ja: B möchte mit G fahren, G aber nicht mit B.}\\[1.5mm]
\textbf{c)} Wer sollte auf jeden Fall nebeneinandersitzen?\\
\lsg[0.98\linewidth]{nur Björn und Doris: Sie haben sich gegenseitig gewählt.}
\end{minipage}
\end{aufgabe}

\begin{merke}[Die Adjazenzmatrix]\raggedright
Ein Graph wird im Computer meist als \textbf{Adjazenzmatrix} (adjazent = benachbart) gespeichert, einer Tabelle wie oben. Zeilen und Spalten stehen für die \lsg[2.2cm]{Knoten}.
\begin{itemize}
\item Führt eine Kante von Knoten $i$ nach Knoten $j$, steht in \lsg[1.9cm]{Zeile $i$} und \lsg[1.9cm]{Spalte $j$} ein X (in Java: \code{true}).
\item Bei gewichteten Graphen steht dort statt des X das \lsg[2.2cm]{Gewicht} der Kante.
\item Bei ungerichteten Graphen ist die Matrix \lsg[2.6cm]{symmetrisch} zur Hauptdiagonalen: Jede Kante $\{i, j\}$ steht zweimal in der Matrix, \lsg[3.8cm]{bei $(i, j)$ und bei $(j, i)$}.
\end{itemize}
\end{merke}

\begin{aufgabe}{2}
\textbf{a)} Gib zu $G_1$ und $G_2$ die Adjazenzmatrix an.\quad
\textbf{b)} Vergleiche die Matrizen. Was fällt bei $G_2$ auf?\\[1mm]
\begin{tabular}{@{}c@{\hspace{4mm}}c@{\hspace{9mm}}c@{\hspace{4mm}}c@{}}
$G_1$\ \begin{tikzpicture}[baseline=(c.base),x=1cm,y=1cm]
\node[kknoten] (a) at (0,1.5) {A}; \node[kknoten] (b) at (1.8,1.5) {B};
\node[kknoten] (c) at (0.9,0.5) {C}; \node[kknoten] (d) at (2.7,0.3) {D};
\draw[gkante] (a)--(b); \draw[gkante] (a)--(c); \draw[gkante] (c)--(b);
\draw[gkante] (b) to[bend left=15] (d); \draw[gkante] (d) to[bend left=15] (b);
\end{tikzpicture}
&
\adjmx[1/2/X,1/3/X,2/4/X,3/2/X,4/2/X]{4}{A,B,C,D}{6.5mm}
&
$G_2$\ \begin{tikzpicture}[baseline=(k3.base),x=1cm,y=1cm]
\node[kknoten] (k1) at (1.2,1.8) {K1}; \node[kknoten] (k2) at (0,0.9) {K2};
\node[kknoten] (k3) at (1.2,0) {K3}; \node[kknoten] (k4) at (2.5,0.5) {K4};
\draw[kante] (k1)-- node[gewicht] {4} (k2); \draw[kante] (k1)-- node[gewicht] {7} (k3);
\draw[kante] (k1)-- node[gewicht] {3} (k4); \draw[kante] (k2)-- node[gewicht] {10} (k3);
\draw[kante] (k3)-- node[gewicht] {5} (k4);
\end{tikzpicture}
&
\adjmx[1/2/4,1/3/7,1/4/3,2/1/4,2/3/10,3/1/7,3/2/10,3/4/5,4/1/3,4/3/5]{4}{K1,K2,K3,K4}{7mm}
\end{tabular}\\[1mm]
\textbf{b)} \lsg[15.5cm]{$G_1$: X-Einträge; $G_2$: Gewichte, Matrix symmetrisch zur Hauptdiagonalen, da ungerichtet.}
\end{aufgabe}

\newpage
\textbf{Implementierung in Java} \web{Projekt Graph}\\[1mm]
\begin{minipage}[t]{0.57\linewidth}
\vspace{0pt}
\begin{java}[basicstyle=\ttfamily\footnotesize]
public class Graph {
   Knoten[] knoten;     // §\lsg[3.8cm]{Knotenfeld}§
   boolean[][] adja;    // §\lsg[3.8cm]{Adjazenzmatrix}§
   int anzKnoten = 5;

   Graph() {
      knoten = new Knoten[§\lsg[2cm]{anzKnoten}§];
      adja = new boolean[§\lsg[2cm]{anzKnoten}§][§\lsg[2cm]{anzKnoten}§];
      // ... Knoten erzeugen
   }
}
\end{java}
\end{minipage}\hfill
\begin{minipage}[t]{0.4\linewidth}
\vspace{0pt}\small\raggedright
Die Matrix ist ein Feld von Feldern (jede Zeile ein Feld): ein \lsg[0.98\linewidth]{zweidimensionales Array}\\[1mm]
Die Knoten werden über ihre \lsg[3cm]{Nummer (Index)} angesprochen. Die Knoten-Objekte liegen im Knotenfeld an derselben Position. \code{adja} muss nicht befüllt werden, \code{boolean} ist anfangs \code{false}.\\[1.5mm]
\textbf{Aufgabe 3)} Starte das Projekt auf der Website unverändert. Was meldet die IDE, und warum?\\
\lsg[0.98\linewidth]{Fehler: knoten ist null}\\[1mm]
Ergänze dann TODO 1 wie im Code links.
\end{minipage}



\begin{aufgabe}[Projekt Graph]{4}
Übertrage die Methode ins Projekt (TODO 2). Sage vorher, was bei den Aufrufen passiert, und teste.\\[1mm]
\begin{minipage}[t]{0.49\linewidth}
\vspace{0pt}
\begin{java}[basicstyle=\ttfamily\scriptsize]
void kanteHinzufuegen(int start, int ende) {
   if (start > anzKnoten || ende > anzKnoten) {
      println("Einfügen nicht möglich.");
   } else {
      adja[start][ende] = true;
      zeichneKante(start, ende);
   }
   printAdja();
}
\end{java}
\end{minipage}\hfill
\begin{minipage}[t]{0.49\linewidth}
\vspace{1mm}\scriptsize\renewcommand{\arraystretch}{1.2}
\begin{tabularx}{\linewidth}{|l|X|}\hline
\textbf{Aufruf} & \textbf{Was passiert?}\\\hline
\code{(0, 1)} & \lsgoder{\lsgfarbe adja[0][1] = true, Kante 0 $\to$ 1}{}\\\hline
\code{(4, 2)} & \lsgoder{\lsgfarbe adja[4][2] = true, Kante 4 $\to$ 2}{}\\\hline
\code{(7, 1)} & \lsgoder{\lsgfarbe 7 > 5: „Einfügen nicht möglich.“}{}\\\hline
\code{(5, 0)} & \lsgoder{\lsgfarbe 5 > 5 falsch: adja[5][0] $\Rightarrow$ Absturz, Index 5 gibt es nicht}{}\\\hline
\end{tabularx}
\end{minipage}\\[1mm]
\textbf{b)} Für welche Art von Graphen passt die Methode? \lsg[5cm]{für gerichtete Graphen}\\[1mm]
\textbf{c)} Verbessere sie: keine ungültigen Nummern, ungerichtete Kanten. Notiere nur die geänderten Zeilen.
\lsglinien{2}{\ttfamily\small if (start < 0 || start >= anzKnoten || ende < 0 || ende >= anzKnoten) \{\\
adja[start][ende] = true;\ \ adja[ende][start] = true;}
\end{aufgabe}

\begin{aufgabe}[Projekt Graph]{5 und 6}
\textbf{5)} Ergänze \code{printAdja()} (TODO 3): Die Matrix soll zeilenweise mit X und O erscheinen.\\
\textbf{6)} \code{nachbarnAusgeben(int k)} gibt alle Knoten aus, zu denen von \code{k} eine Kante führt (TODO 4).\\[1mm]
\begin{minipage}[t]{0.58\linewidth}
\vspace{0pt}
\begin{java}[basicstyle=\ttfamily\scriptsize]
void printAdja() {
   for (int zeile = 0; zeile < §\lsg[1.5cm]{anzKnoten}§; zeile++) {
      for (int spalte = 0; spalte < §\lsg[1.4cm]{anzKnoten}§; §\lsg[1.1cm]{spalte++}§) {
         if (adja[§\lsg[0.9cm]{zeile}§][§\lsg[0.9cm]{spalte}§]) {
            print("X ");
         } else {
            print("O ");
         }
      }
      §\lsg[1.8cm]{println();}§ // Zeilenumbruch
   }
   println();
}
\end{java}
\end{minipage}\hfill
\begin{minipage}[t]{0.4\linewidth}
\vspace{0pt}
\begin{java}[basicstyle=\ttfamily\scriptsize]
void nachbarnAusgeben(int k) {
   print("Nachbarn von " + k + ": ");
   for (int i = 0; i < §\lsg[1.4cm]{anzKnoten}§; i++) {
      if (§\lsg[2cm]{adja[k][i]}§) {
         print(§\lsg[1.6cm]{i + " "}§);
      }
   }
   println();
}
\end{java}
\end{minipage}\\[1mm]
\textbf{Merke:} Die Nachbarn von Knoten $k$ stehen in \lsg[2cm]{Zeile $k$}. Man durchläuft sie mit einer \lsg[5cm]{Wiederholung mit fester Anzahl} und prüft jeden Eintrag.
\end{aufgabe}


\end{document}
